\chapter{Resultados R7 y estado de la evidencia}

\COBRIAAlcanceEPocho

\section{Alcance histórico y evaluación posterior}

% Generado por editorial/ep6/build.mjs; no editar.
\begin{table}[htbp]\centering\small
\caption{Estratos de evidencia y límites de interpretación.}
\begin{tabularx}{\textwidth}{p{2.1cm}X X}\toprule
Estrato & Aporte disponible & Inferencia no establecida\\\midrule
Norma & Contrato e invariantes de Core 1.0 & Eficacia o conformidad por el solo uso del nombre\\
Observación & Dos inventarios originales y un perfil complementario & Prevalencia, causalidad o muestra independiente\\
R7 histórico & 270 corridas; ocho celdas favorables y una no concluyente & Validación retrospectiva de PC2 o superioridad del contrato\\
Microensayo & 630 observaciones en un proceso & Rendimiento distribuido de arquitecturas completas\\
PC2 local & Pruebas de persistencia, publicación, autorización y consumidor & Nuevas cifras en tres motores o calidad de un LLM\\
PC3 pendiente & Diseño y análisis preparado sin observaciones nuevas & Resultado confirmatorio, celdas imputadas o mejora estimada\\
Evaluación externa & Protocolos y material de transferencia & Revisión independiente, aprendizaje o productividad demostrados\\
\bottomrule\end{tabularx}\end{table}


La propuesta combina mecanismos existentes, no los sustituye por una
afirmación de novedad absoluta. Separar lectura y escritura puede añadir complejidad
sin beneficio cuando el dominio no lo justifica \parencite{fowler2011cqrs}.
Reconstruir una proyección desde documentos canónicos tampoco exige conservar todos
los cambios como eventos \parencite{fowler2005eventsourcing}. El control de revisión
de un documento debe distinguirse de una transacción entre documentos
\parencite{apacheCouchDocumentApi}. Estas referencias orientan alternativas y amenazas;
no demuestran superioridad de COBRIA. La selección bibliográfica exploratoria asistida
no equivale al doble cribado humano ni a una revisión sistemática completa.


R7 conserva los resultados de su arnés original. La aprobación de sus controles no
equivale a aprobar requisitos incorporados después ni demuestra comprensión por
personas usuarias. PC2 congela contratos adicionales de persistencia, publicación,
autorización previa al almacenamiento y abstención de un consumidor determinista.
Sus pruebas locales no son una réplica confirmatoria de tres motores ni una evaluación
de modelos generativos. PC3 permanece sin observaciones nuevas mientras no se ejecute
el diseño congelado con conectores y motores verificados.

La publicación mediante comparación e intercambio de un agregado no establece
atomicidad entre documentos independientes. El filtrado por ámbito comprueba exclusión
de contenido en las rutas ensayadas, no ausencia de canales laterales temporales.
La revisión independiente, la réplica externa y el estudio de transferencia mantienen
su identidad separada de las pruebas automáticas del autor.

Este capítulo interpreta el conjunto histórico R7 sin mezclarlo con controles
posteriores o mediciones preliminares. La fuente del análisis es
\path{replication/results/core-r7-statistics/summary.json}; sus entradas son los
eventos crudos cuyos SHA-256 se registran por motor. El alcance es interno: un nodo por
motor y configuraciones declaradas. No es una certificación ni una comparación universal.

\section{Diseño y conformidad}

R7 reunió \COBRIARSevenObservations{} corridas internas, distribuidas en partes iguales
entre MongoDB \COBRIARSevenMongoVersion, CouchDB \COBRIARSevenCouchVersion{} y
OpenSearch \COBRIARSevenOpenSearchVersion. Cada motor se evaluó con tamaños 100, 1 000
y 5 000, dos calentamientos y \COBRIARSevenRunsPerCell{} corridas válidas por celda en
orden aleatorizado y reproducible. Las 90 corridas de cada motor aprobaron los
indicadores P1, P2, R1, S1 y A1 del arnés histórico, no la totalidad de los controles PC2.

\COBRIALimiteHistoricoRSeven

\begin{table}[H]\centering\small
\caption{Resultados pareados de lectura de R7.}
\begin{tabular}{lrrrrr}\toprule
Motor & $n$ & razón mediana & diferencia media (ms) & $d_z$ & $p$ de Holm\\\midrule
\COBRIARSevenTableRows
\bottomrule\end{tabular}
\end{table}

Una razón menor que uno favorece a la proyección. Ocho celdas fueron concluyentes y
favorables después del ajuste de Holm. OpenSearch con $n=100$ tuvo diferencia media de
0,020 ms, $d_z=0,002$ y $p_{Holm}=0,993$; por tanto se informa como no concluyente.

\section{Alcance de la inferencia}

Los resultados describen comprobaciones del arnés y costos en los entornos capturados.
No demuestran conformidad normativa integral, superioridad universal, seguridad exhaustiva, productividad
humana ni desempeño multinodo. La microcomparación de siete estrategias contiene
\COBRIARSevenAlternatives{} observaciones, pero se ejecuta en un proceso Node.js y solo
sirve como contraste algorítmico. Permanecen pendientes la réplica externa independiente,
el estudio de transferencia y el doble cribado bibliográfico.

\section{Análisis histórico y diseño posterior}

\COBRIAAnalisisSeparado
